\chapter{Establishing API Communication}

This chapter connects Chapters 1--7 into one full round trip: a React page
calling \texttt{api.get("/tasks")} and rendering what MongoDB returns. Every
CRUD feature later in the book is a variation of this pattern.

\section{The Complete Cycle, Labeled}

\begin{diagrambox}[GET /api/tasks -- Full Cycle]
\begin{verbatim}
Dashboard.jsx: useEffect(() => taskService.getTasks(), [])
  -> taskService.js: api.get("/tasks")
  -> axios instance: baseURL + withCredentials (JWT cookie)
  -> HTTP Request: GET /api/tasks, Cookie: token=<jwt>
  -> app.js: app.use("/api/tasks", taskRoutes)
  -> taskRoutes.js: router.get("/", protect, getTasks)
  -> middleware/auth.js: protect verifies JWT -> req.user
  -> taskController.js: getTasks(req, res)
  -> Mongoose: Task.find({ user: req.user._id })
  -> MongoDB: scans "tasks", returns matching documents
  -> taskController.js: res.status(200).json({ success:true, data:tasks })
  -> HTTP Response: 200 OK, { success, data: [...] }
  -> axios resolves with response.data
  -> taskService.js returns res.data.data
  -> Dashboard.jsx: setTasks(data) -> TaskList renders TaskCards
\end{verbatim}
\end{diagrambox}

\section{The Frontend Side}

\begin{lstlisting}[language=JavaScript]
// src/services/taskService.js
import api from "../api/axios";

export async function getTasks() {
  try {
    const res = await api.get("/tasks");
    return res.data.data;
  } catch (err) {
    throw new Error(err.response?.data?.message || "Failed to load tasks");
  }
}
\end{lstlisting}

\begin{lstlisting}[language=JavaScript]
// src/pages/Dashboard.jsx
import { useEffect, useState } from "react";
import { getTasks } from "../services/taskService";
import TaskList from "../components/TaskList";

export default function Dashboard() {
  const [tasks, setTasks] = useState([]);
  const [error, setError] = useState("");

  useEffect(() => {
    getTasks()
      .then(setTasks)
      .catch((err) => setError(err.message));
  }, []);

  if (error) return <p className="error">{error}</p>;
  return <TaskList tasks={tasks} />;
}
\end{lstlisting}

\section{The Backend Side}

\begin{lstlisting}[language=JavaScript]
// routes/taskRoutes.js
const express = require("express");
const router = express.Router();
const { protect } = require("../middleware/auth");
const { getTasks, createTask } = require("../controllers/taskController");

router.get("/", protect, getTasks);
router.post("/", protect, createTask);

module.exports = router;
\end{lstlisting}

\begin{lstlisting}[language=JavaScript]
// controllers/taskController.js
const Task = require("../models/Task");

exports.getTasks = async (req, res) => {
  try {
    const tasks = await Task.find({ user: req.user._id }).sort({ createdAt: -1 });
    res.status(200).json({ success: true, data: tasks });
  } catch (err) {
    res.status(500).json({ success: false, message: "Failed to fetch tasks" });
  }
};
\end{lstlisting}

\begin{lstlisting}[language=JavaScript]
// app.js (wiring)
const express = require("express");
const app = express();

app.use(express.json());
app.use("/api/tasks", require("./routes/taskRoutes"));
app.use("/api/auth", require("./routes/authRoutes"));

module.exports = app;
\end{lstlisting}

\begin{notebox}[NOTE]
Each layer only knows the layer next to it -- the controller never touches
\texttt{req}/\texttt{res} inside Mongoose calls, and routing never touches MongoDB
directly. That's what keeps the chain easy to debug and test.
\end{notebox}

\section{Standard Response Shapes}

Every endpoint in this book returns one of these shapes, keeping the
frontend's \texttt{res.data.data} unwrapping consistent.

\begin{table}[h]
\centering
\begin{tabularx}{\textwidth}{lXl}
\toprule
\textbf{Shape} & \textbf{Used for} & \textbf{Example} \\
\midrule
\texttt{\{ success, data \}} & Fetch/create/update & \texttt{GET/POST /tasks} \\
\texttt{\{ success, message \}} & No payload actions & \texttt{DELETE /tasks/:id} \\
\texttt{\{ success, data, page, limit, total, totalPages \}} & Pagination (Ch. 20) & \texttt{GET /tasks?page=2} \\
\bottomrule
\end{tabularx}
\end{table}

\begin{warnbox}[WARNING]
Never mix shapes for the same endpoint across success/error cases -- keep
\texttt{success: boolean} as the first field always.
\end{warnbox}

\begin{tipbox}[TIP]
When a feature breaks, check in order: URL, HTTP method, route match,
controller query field, response shape. This covers most MERN bugs.
\end{tipbox}
